% Part: computability
% Chapter: recursive-functions
% Section: primes

\documentclass[../../../include/open-logic-section]{subfiles}

\begin{document}

\olfileid{cmp}{rec}{pri}
\olsection{અવિભાજ્ય સંખ્યાઓ}

સીમિત પરિમાણીકરણ અને સીમાબદ્ધ લઘુત્તમીકરણથી ઘણા સ્વાભાવિક વિધેયો
અને સંબંધો આદિમ પુનરાવર્તી છે તે બતાવવા માટે પૂરતાં સાધનો મળે છે.
દાખલા તરીકે, ``$x$ એ $y$નો અવયવ છે'' એવો સંબંધ વિચારો, જેને
$x \mid y$ લખીએ છીએ. $y$ને~$x$થી ભાગતાં શેષ ન રહે, એટલે $y$ એ
~$x$નો પૂર્ણાંક ગુણિત હોય, ત્યારે $x \mid y$ લાગુ પડે છે. (ભાગાકારક
શૂન્ય ન હોય અને $y$ને $x$થી ભાગતાં શેષ $> 0$ હોય, તો $x \nmid y$
લખીએ છીએ. શૂન્ય ભાગાકારક માટે નીચેનું સમીકરણ સંબંધ નક્કી કરે છે.)
બીજા શબ્દોમાં, કોઈ~$z$ માટે $x \cdot z = y$ હોય
ત્યારે અને માત્ર ત્યારે જ $x \mid y$ છે. એવો $z$ અસ્તિત્વમાં
હોય, તો કોઈ એવો $z$ $\leq y$ પસંદ કરી શકાય છે. તેથી કોઈ
$z \le y$ માટે $x \cdot z = y$ હોય ત્યારે અને માત્ર ત્યારે જ
$x \mid y$ છે. સમાનતા~$=$ અને ગુણાકારમાંથી સીમિત
અસ્તિત્વપરિમાણીકરણ વડે $x \mid y$ સંબંધને આમ વ્યાખ્યાયિત કરી
શકીએ:
\[
x \mid y \defiff \bexists{z \leq y}{(x \cdot z) = y}.
\]
આથી $x \mid y$ આદિમ પુનરાવર્તી છે એવું આપણે બતાવ્યું છે.

\sourcecorrection{OLCMP-005}{સ્થિર અંગ્રેજી મૂળના કૌંસમાં
``$x$ને $y$થી ભાગતાં'' શેષ એવું લખાયું છે. $x \mid y$ના પહેલાંના
વાક્ય અને વ્યાખ્યાકારક સમીકરણ પ્રમાણે ભાગવાનું $y$ને $x$થી છે.
શૂન્ય ભાગાકારક માટે ભાગાકારનો શેષ વ્યાખ્યાયિત નથી; ગુજરાતી કૌંસમાં
આ અપવાદ પણ સ્પષ્ટ કર્યો છે.}
\sourcecorrection{OLCMP-006}{સ્થિર અંગ્રેજી મૂળ કહે છે કે
$x \cdot z = y$ સંતોષતો ``કોઈ પણ'' $z$ જરૂર $\leq y$ છે. $x=y=0$
હોય ત્યારે મોટા $z$ પણ સમીકરણ સંતોષે છે. સીમિત પરિમાણીકરણ માટે
માત્ર એવો ઓછામાં ઓછો એક $z \leq y$ મળે એટલું જરૂરી અને સાચું છે;
ગુજરાતી વાક્યમાં આ દાવો રાખ્યો છે.}

પ્રાકૃતિક સંખ્યા~$x$ \emph{અવિભાજ્ય} કહેવાય જો તે $0$ કે $1$
ન હોય અને માત્ર $1$ તથા પોતાનાથી જ વિભાજ્ય હોય. બીજા શબ્દોમાં,
અવિભાજ્ય સંખ્યાઓ માટે $y \mid x$ હોય ત્યારે $y = 1$ અથવા~$y=x$
હોય છે. $x$ અવિભાજ્ય છે કે નહીં તે તપાસવા $y \le x$ માટે જ
$y \mid x$ તપાસવું પૂરતું છે, કારણ કે $y > x$ હોય તો આપોઆપ
$y \nmid x$. તેથી $x$ અવિભાજ્ય હોય ત્યારે અને માત્ર ત્યારે જ
લાગુ પડતો સંબંધ $\fn{Prime}(x)$ આમ વ્યાખ્યાયિત કરી શકાય:
\[
\fn{Prime}(x) \defiff x \geq 2 \land \bforall{y \leq x}{(y \mid x \lif y
  = 1 \lor y = x)}
\]
અને આથી તે આદિમ પુનરાવર્તી છે.

અવિભાજ્ય સંખ્યાઓ $2$, $3$, $5$, $7$, $11$ વગેરે છે. આ શ્રેણીમાં
$x$મી અવિભાજ્ય સંખ્યા આપતું વિધેય $p(x)$ વિચારો: $p(0) =
2$, $p(1) = 3$, $p(2) = 5$ વગેરે. (સગવડ માટે આપણે ઘણી વાર
$p(x)$ને $p_x$ પણ લખીશું: $p_0=2$, $p_1=3$ વગેરે.)

$x$થી મોટી પહેલી અવિભાજ્ય સંખ્યા આપતું વિધેય
$\fn{nextPrime(x)}$ હોય, તો $p$ને આદિમ પુનરાવર્તન વડે સરળતાથી
વ્યાખ્યાયિત કરી શકાય:
\begin{align*}
  p(0) & = 2\\
  p(x+1) & = \fn{nextPrime}(p(x))
\end{align*}
$\fn{nextPrime}(x)$ એ $y > x$ અને $y$ અવિભાજ્ય હોય એવી સૌથી
નાની $y$ છે; અનિયત મર્યાદા સુધીની શોધ વડે તેને સરળતાથી ગણી
શકાય છે. પરંતુ યુક્લિડના પરિણામને કારણે તેને સીમાબદ્ધ
લઘુત્તમીકરણ વડે પણ વ્યાખ્યાયિત કરી શકાય: $x$ અને $\fact{x}+1$
વચ્ચે હંમેશાં એક અવિભાજ્ય સંખ્યા મળે છે.
\[
  \fn{nextPrime(x)} =
  \bmin{y \leq \fact{x}+1}{(y > x \land \fn{Prime}(y))}.
\]
આથી $\fn{nextPrime}(x)$ અને પરિણામે $p(x)$ માત્ર સંગણનીય જ નહીં,
પણ આદિમ પુનરાવર્તી છે.

(જિજ્ઞાસા હોય તો યુક્લિડના પરિણામની ટૂંકી સાબિતી આ છે. ધારો કે
$p_n$ એ $\le x$ એવી સૌથી મોટી અવિભાજ્ય સંખ્યા છે; અને
$\le x$ એવી બધી અવિભાજ્ય સંખ્યાઓનો ગુણાકાર $p =
p_0\cdot p_1 \cdot \dots \cdot p_n$ વિચારો. $p+1$ પોતે અવિભાજ્ય
હોય અથવા $x$ અને~$p+1$ વચ્ચે કોઈ અવિભાજ્ય સંખ્યા હોય. શા માટે?
ધારો કે $p+1$ અવિભાજ્ય નથી. ત્યારે કોઈ અવિભાજ્ય સંખ્યા
$q \mid p+1$ છે, જ્યાં $q < p+1$. $\le x$ એવી એકેય અવિભાજ્ય
સંખ્યા $p+1$ને ભાગતી નથી. ($p$ની વ્યાખ્યા પ્રમાણે દરેક
$p_i \le x$ એ~$p$ને ભાગે છે, એટલે શેષ~$0$ છે. તેથી દરેક
$p_i \le x$ વડે $p+1$ને ભાગતાં શેષ~$1$ રહે છે; એટલે
$p_i \nmid p+1$.) આથી $q$ એ $> x$ અને $< p+1$ એવી અવિભાજ્ય
સંખ્યા છે. તથા $p \le \fact{x}$ હોવાથી $> x$ અને
$\le \fact{x}+1$ એવી અવિભાજ્ય સંખ્યા મળે છે.)

\begin{prob}
સીમાબદ્ધ લઘુત્તમીકરણ વડે પૂર્ણાંક ભાગાકાર $d(x, y)$ વ્યાખ્યાયિત
કરો.
\end{prob}

\end{document}
